%% This file contains a number of tweaks that are typically applied to the main document.
%% They are not enabled by default, but can be enabled by uncommenting the relevant lines.


%%
%% Official support for 2-column teaser figure
%%
\usepackage{cuted} %< for strip environment

%%
%% For auto-referencing of figures (see fig/teaser.tex and \label{} therein) 
%%
\usepackage{currfile} %< for \currfilebase macro

%%
%% caption and (global) spacing overrides
%%
\usepackage{caption} %< for "\captionof" commands (teaser)
% \setlength{\abovecaptionskip}{.5em} % space above caption
% \setlength{\belowcaptionskip}{1em} % space below caption

%%
%% Inline annotations; for predefined colors, refer to "dvipsnames" in the xcolor package:
%% https://tinyurl.com/overleaf-colors
%%
\newcommand{\red}[1]{{\color{red}#1}}
\newcommand{\todo}[1]{{\color{red}#1}}
\newcommand{\TODO}[1]{\textbf{\color{red}[TODO: #1]}}
%%
%% Disable for camera ready / submission by uncommenting these lines:
%%
% \renewcommand{\TODO}[1]{}
% \renewcommand{\todo}[1]{#1}

%%
%% Work harder in optimizing text layout. Typically shrinks text by 1/6 of page, enable
%% it at the very end of the writing process, when you are just above the page limit.
%%
% \usepackage{microtype}

%%
%% Fine-tune paragraph spacing.
%%
% \renewcommand{\paragraph}[1]{\vspace{.5em}\noindent\textbf{#1.}}

%%
%% Globally adjusts space between figure and caption.
%%
% \setlength{\abovecaptionskip}{.5em}

%%
%% Allows "the use of \paper to refer to the project name"
%% with automatic management of space at the end of the word.
%%
% \usepackage{xspace}
% \newcommand{\paper}{ProjectName\xspace}

%%
%% Commonly used math definitions.
%%
% \DeclareMathOperator*{\argmin}{arg\,min}
% \DeclareMathOperator*{\argmax}{arg\,max}

%%
%% Tighten underline.
%%
% \usepackage{soul}
% \setuldepth{foobar}

%%
%% Reconfigure paragraph (tweaks spacing)
%%

%%
%% Project additions (TALON), loaded before hyperref as the kit asks.
%%
% Workaround for lineno v5.7 (TeX Live 2026) with LaTeX >= 2025-06: its new column hook tags the wrong column,
% so right-column line numbers land in the gutter on top of the text. Restore lineno's pre-2025 column tagging.
\makeatletter
\@ifpackagelater{lineno}{2025/06/01}{%
  \RemoveFromHook{build/column/before}[lineno]%
  \AtBeginDocument{%
    \def\@LN@makecol{%
      \@LN@orig@makecol
      \setbox\@outputbox\vbox{%
        \boxmaxdepth\@maxdepth
        \protected@write\@auxout{}{\string\def\string\@LN@column{\if@firstcolumn1\else2\fi}}%
        \box\@outputbox}}}%
}{}
\makeatother

\usepackage{amsthm}
\usepackage{tikz}
\usepackage{algorithm}
\usepackage[noend]{algpseudocode}   % no "end for"/"end if" lines (DDPM / CLRS style)
\algrenewcommand\algorithmicrequire{\textbf{Input:}}
\algrenewcommand\algorithmicensure{\textbf{Output:}}
\algrenewcommand\algorithmicindent{1em}
\usetikzlibrary{arrows.meta, positioning, fit, backgrounds, calc, shapes.geometric, shapes.callouts}

% Propositions: plain run-in heads without a frame and without LaTeX list machinery. A tcolorbox frame made each
% proposition an unbreakable block, which the kit's \flushbottom turned into large white gaps. A list-based
% theorem (amsthm or \newtheorem) can lose its head into the page corner under the review-mode line numbering.
\newcounter{proposition}
\NewDocumentEnvironment{proposition}{o}
  {\par\addvspace{\topsep}\refstepcounter{proposition}\noindent
   \textbf{Proposition~\theproposition}\IfValueT{#1}{ (#1)}\textbf{.}\ \itshape\ignorespaces}
  {\par\addvspace{\topsep}}
\newtheorem{lemma}{Lemma}
\newcommand{\name}{TALON\xspace}
\newcommand{\one}{\mathbf{1}}
\DeclareMathOperator*{\argmax}{arg\,max}
\DeclareMathOperator*{\argmin}{arg\,min}

% TikZ colours used by the overview figure (fig/overview.tex)
\definecolor{orxblue}{HTML}{0072B2}
\definecolor{orxorange}{HTML}{E69F00}
\definecolor{orxgreen}{HTML}{009E73}
\definecolor{orxred}{HTML}{D55E00}
\definecolor{orxpurple}{HTML}{CC79A7}
\definecolor{orxsky}{HTML}{56B4E9}

\newcommand{\ayush}[1]{\textcolor{purple}{[Ayush: #1]}}
\newcommand{\fan}[1]{\textcolor{blue}{[Fan: #1]}}
\newcommand{\francis}[1]{\textcolor{orange}{[Francis: #1]}}